Intuitively, \mmt is a declarative language for theories and views over an arbitrary object language.
Its treatment of object languages is abstract enough to subsume most logics and type theories which are practically relevant.

Fig.~\ref{fig:mmt} gives an overview of the fundamental MMT concepts.
In the simplest case, \textbf{theories} $\Sigma$ are lists of \textbf{constant declarations} $c:E$, where $E$ is an expression that may use the previously declared constants.
Naturally, $E$ must be subject to some type system (which MMT is also parametric in), but the details of this are not critical for our purposes here.
We say that $\Sigma'$ includes $\Sigma$ if it contains every constant declaration of $\Sigma$.

\begin{figure}[htb]\centering
\begin{tabular}{|l|l|l|}\hline
\multicolumn{3}{|c|}{meta-theory: a fixed theory $M$}\\\hline
       & Theory $\Sigma$    & View $\sigma:\Sigma\to\Sigma'$ \\\hline
set of & typed constant declarations $c:E$ & assignments $c\mapsto E'$ \\
$\Sigma$-expressions $E$ & formed from $M$- and $\Sigma$-constants  & mapped to $\Sigma'$ expressions \\\hline
\end{tabular}
\caption{Overview of MMT Concepts}\label{fig:mmt}
\end{figure}

Correspondingly, a \textbf{view} $\sigma:\Sigma\to\Sigma'$ is a list of \textbf{assignments} $c\mapsto e'$ of $\Sigma'$-expressions $e'$ to $\Sigma$-constants $c$.
To be well-typed, $\sigma$ must preserve typing, i.e., we must have $\vdash_{\Sigma'}e':\ov{\sigma}(E)$.
Here $\ov{\sigma}$ is the homomorphic extension of $\sigma$, i.e., the map of $\Sigma$-expressions to $\Sigma'$-expressions that substitutes every occurrence of a $\Sigma'$-constant with the $\Sigma'$-expression assigned by $\sigma$.
We call $\sigma$ \textbf{simple} if the expressions $e'$ are always $\Sigma'$-\emph{constants} rather than complex expressions.
The type-preservation condition for an assignment $c\mapsto c'$ reduces to $\ov{\sigma}(E)=E'$ where $E$ and $E'$ are the types of $c$ and $c'$.
We call $\sigma$ \textbf{partial} if it does not contain an assignment for every $\Sigma$-constant and \textbf{total} otherwise.
A partial view from $\Sigma$ to $\Sigma'$ is the same as a total view from some theory included by $\Sigma$ to $\Sigma'$.

Importantly, we can then show generally at the MMT-level that if $\sigma$ is well-typed, then $\ov{\sigma}$ preserves all typing and equality judgments over $\Sigma$.
In particular, if we represent proofs as typed terms, views preserve the theoremhood of propositions.
This property makes views so valuable for structuring, refactoring, and integrating large corpora.

MMT achieves language-independence through the use of \textbf{meta-theories}: every MMT-theory may designate a previously defined theory as its meta-theory.
For example, when we represent the HOL Light library in MMT, we first write a theory $L$ for the logical primitives of HOL Light.
Then each theory in the HOL Light library is represented as a theory with $L$ as its meta-theory.
In fact, we usually go one step further: $L$ itself is a theory, whose meta-theory is a logical framework such as LF.
That allows $L$ to concisely define the syntax and inference system of HOL Light.

However, for our purposes, it suffices to say that the meta-theory is some fixed theory relative to which all concepts are defined.
Thus, we assume that $\Sigma$ and $\Sigma'$ have the same meta-theory $M$, and that $\ov{\sigma}$ maps all $M$-constants to themselves.


\begin{wrapfigure}r{4.2cm}\vspace*{-2em}
\begin{tabular}{l@{\quad}c@{\quad}l}
 $\Gamma$	& $::=$	& $(x:E)^\ast$ \\
 $E$      & $::=$ & $c \mid x \mid \ombind{E}{\Gamma}{E^+}$ \\
\end{tabular}\vspace*{-2em}
\end{wrapfigure}
It remains to define the exact syntax of expressions.   
In the grammar on the right $c$ refers to constants (of the meta-theory or previously declared in the current theory) and $x$ refers to bound variables.
Complex expressions are of the form $\ombind{o}{x_1:t_1,\ldots,x_m:t_m}{a_1,\ldots,a_n}$, where
\begin{compactitem}
 \item $o$ is the operator that forms the complex expression,
 \item $x_i:t_i$ declares variable of type $t_i$ that are bound by $o$ in subsequent variable declarations and in the arguments,
 \item $a_i$ is an argument of $o$
\end{compactitem}
The bound variable context may be empty, and we write $\oma{o}{\vec{a}}$ instead of $\ombind{o}{\cdot}{\vec{a}}$.
For example, the axiom $\forall x:\cn{set},y:\cn{set}.\; \cn{beautiful}(x) \wedge y
\subseteq x \Rightarrow \cn{beautiful}(y)$ would instead be written as \[\ombind{\forall}{x:\cn{set},y:\cn{set}}{\oma{\Rightarrow}{\oma{\wedge}{\oma{\cn{beautiful}}{x},\oma{\subseteq}{y,x}},\oma{\cn{beautiful}}{y}}}\]

Finally, we remark on a few additional features of the MMT language that are important for large-scale case studies but not critical to understand the basic intuitions of results.
MMT provides a module system that allows theories to instantiate and import each other. The module system is conservative: every theory can be \emph{elaborated} into one that only declares constants.
MMT constants may carry an optional definiens, in which case we write $c:E=e$.
Defined constants can be eliminated by definition expansion.


%\begin{example}[$\lambda$-calculus]\label{ex:lamsyn}
%To give an urtheory for the simply-typed $\lambda$-calculus in \mmt, we use $4$ constructors $\mathcal{C}=\{\type,\to,\lambda,@\}$.
%Alternatively, we can add constructors $\kind$ and $\Pi$ to obtain an urtheory for the dependently-typed $\lambda$-calculus.
%
%We introduce the usual notations for them as follows:
%
%\begin{center}
%\begin{tabular}{|l|l|l|}
%\hline
%Constructor       & Abstract \mmt expression     & Concrete notation\\
%\hline
%universe of types & $\cons{\type}{\cdot}{\cdot}$ & $\type$ \\
%function types    & $\cons{\to}{\cdot}{A,B}$     & $A\to B$  \\
%abstraction       & $\cons{\lambda}{x:A}{t}$     & $\lambda x:A.t$ \\
%application       & $\cons{@}{\cdot}{f,t}$       & $f\,t$\\
%\hline
%universe of kinds & $\cons{\kind}{\cdot}{\cdot}$ & $\kind$ \\
%dependent function types & $\Pi(x:A;t)$  & $\Pi x:A.t$ \\
%\hline
%\end{tabular}
%\end{center}
%\end{example}

%\begin{oldpart}{FR: replaced with the above}
%For the purposes of this paper, we will work with the (only slightly simplified) grammar given in Figure \ref{fig:mmtgrammar}.
%
%\begin{figure}[ht]\centering\vspace*{-1em}
%  \begin{mdframed}
%    \begin{tabular}{rl@{\quad}l|l}
%      \cn{Thy} 	$::=$ 	& $T[:T]=\{ (\cn{Inc})^\ast\ (\cn{Const})^\ast \}$	& Theory				& \multirow{2}{*}{Modules} \\
%      \cn{View} 	$::=$	& $T : T \to T = \{ (\cn{Ass})^\ast \}$			& View 				& \\ \hline
%      \cn{Const} 	$::=$ 	& $C [:t] [=t]$ 								& Constant Declarations& \multirow{3}{*}{Declarations} \\
%      \cn{Inc}	$::=$	& $\incl T$								& Includes			& \\
%      \cn{Ass}	$::=$	& $C = t$									& Assignments		& \\ \hline
%      $\Gamma$	$::=$	& $(x [:t][=t])^\ast$						& Variable Contexts	& \multirow{2}{*}{Objects} \\
%      $t$ 		$::=$ 	& $x \mid T?C \mid \oma{t}{(t)^+} \mid \ombind{t}{\Gamma}{t} $		& Terms 	& \\\hline\hline
%      $T$		$::=$ 	& \cn{String}								& Module Names 		& \multirow{3}{*}{Strings} \\
%      $C$		$::=$ 	& \cn{String}								& Constant Names		& \\
%      $x$		$::=$ 	& \cn{String}								& Variable Names		& \\
%    \end{tabular}
%  \end{mdframed}\vspace*{-.5em}
%  \caption{The MMT Grammar}\label{fig:mmtgrammar}\label{fig:mmtgrammar}\vspace*{-1em}
%\end{figure}
%
%In more detail:
%\begin{itemize}
%	\item Theories have a module name, an optional \emph{meta-theory} and a body consisting of \emph{includes} of other theories
%		 and a list of \emph{constant declarations}.
%	\item Constant declarations have a name and two optional term components; a \emph{type} ($[:t]$), and a \emph{definition} ($[=t]$).
%	\item Views $V : T_1 \to T_2$ have a module name, a domain theory $T_1$, a codomain theory $T_2$ and a body consisting of assignments
%		$C = t$.
%	\item Terms are either 
%		\begin{itemize}
%			\item variables $x$, 
%			\item symbol references $T?C$ (referencing the constant $C$ in theory $T$), 
%			\item applications $\oma{f}{a_1,\ldots,a_n}$ of a term $f$ to a list of arguments $a_1,\ldots,a_n$ or
%			\item binding application $\ombind{f}{x_1[:t_1][=d_1],\ldots,x_n[:t_n][=d_n]}{b}$, where $f$ \emph{binds} the variables
%				$x_1,\ldots,x_n$ in the body $b$ (representing binders such as quantifiers, lambdas, dependent type constructors etc.).
%		\end{itemize}
%\end{itemize}
%The term components of a constant in a theory $T$ may only contain symbol references to constants declared previously in $T$, or that are declared in some theory $T'$ (recursively) included in $T$ (or its meta-theory, which we consider an \emph{include} as well). 
%We can eliminate all includes in a theory $T$ by simply copying over the constant declarations in the included theories; we call this process \emph{flattening}. We will often and without loss of generality assume a theory to be \emph{flat} for convenience.
%
%An assignment in a view $V:T_1\to T_2$ is syntactically well-formed if for any assignment $C=t$ contained, $C$ is a constant declared in the flattened domain $T_1$ and $t$ is a syntactically well-formed term in the codomain $T_2$. We call a view \emph{total} if all \emph{undefined} constants in the domain have a corresponding assignment and \emph{partial} otherwise.
%
%\end{oldpart}

%\ifshort\else
\subsection{Proof Assistant Libraries in MMT}\label{sec:oaf}

As part of the OAF project~\cite{OAFproject:on}, we have imported several proof assistant libraries into the MMT system. To motivate some of the design choices made in this paper, we will outline the general procedure behind these imports.

\paragraph{} First, we formalize the core logical foundation of the system. We do so by using the logical framework LF~\cite{lf} (at its core a dependently-typed lambda calculus) and various extensions thereof, which are implemented in and supported by the MMT system. In LF, we can formalize the foundational primitives using the usual judgments-as-types and higher-order abstract syntax encodings -- hence theorems and axioms are represented as constants with a type $\vdash P$ for some proposition $P$, and primitive constructs like lambdas are formalized as LF-functions taking LF-lambda-expressions -- which serve as a general encoding of any variable binders -- as arguments.

The resulting formalizations are then used as meta-theory for imports of the libraries of the system under consideration. This results in a theory graph as in Figure \ref{fig:oaf}.

\begin{figure}[ht]\centering
 \begin{tikzpicture}
 \node (MMT) at (2,2.5) {MMT};
 
 \draw[fill=orange!40] (2,1) ellipse (1.5cm and .6cm);
 \node[color=orange] at (-3.3,1) {Logical Frameworks};
 \node (L) at (1,1) {LF};
 \node (Lx) at (3,1) {LF+X};
 \draw[arrow](MMT) -- (L);
 \draw[arrow](MMT) -- (Lx);
 \draw[arrow](L) -- (Lx);

 \draw[fill=red!40] (2,-.5) ellipse (3.2cm and .6cm);
 \node[color=red] at (-3.3,-.5) {Foundations};
 \node at (2,-.7) {\ldots};

 \draw[fill=blue!40] (0,-2.25) ellipse (1.9cm and .8cm);

 \node (H) at (0,-.5) {HOL Light};
 \node[color=blue!80] at (-3.5,-2) {HOL Light library};
 \node (B) at (-1,-2) {Bool};
 \node (A) at (1,-2) {Arith};
 \node (E) at (0,-2.5) {\ldots};
 \draw[arrow](L) -- (H);
 \draw[arrow](H) -- (B);
 \draw[arrow](H) -- (A);
 \draw[arrow](B) -- (A);

 \draw[fill=olive!40] (4,-2.25) ellipse (1.9cm and .8cm);

 \node (M) at (4,-.5) {PVS};
 \node[color=olive] at (-3.3,-2.5) {PVS library};
 \node (B') at (3,-2) {Booleans};
 \node (A') at (5,-2) {Reals};
 \node (E') at (4,-2.5) {\ldots};

 \node (A) at (1,-2) {Arith};
 \node (E) at (0,-2.5) {\ldots};
 \draw[arrow](Lx) -- (M);
 \draw[arrow](M) -- (B');
 \draw[arrow](M) -- (A');
 \draw[arrow](B') -- (A');
 \end{tikzpicture}
 \caption{A (Simplified) Theory Graph for the OAF Project}\label{fig:oaf}
\end{figure}
%\fi



%%% Local Variables:
%%% mode: latex
%%% eval:  (set-fill-column 5000)
%%% eval:  (visual-line-mode)
%%% TeX-master: "paper"
%%% End:

%  LocalWords:  fig:mmt textbf textbf mapsto vdash_ emph hline ctabular ldots,x_m:t_m vec
%  LocalWords:  ldots,a_n compactitem cdot vec ednote ex:lamsyn urtheory mathcal f,t y,x
%  LocalWords:  x:A;t oldpart fig:mmtgrammar centering vspace mdframed multirow ldots,x_n
%  LocalWords:  OAFproject:on vdash formalized formalizations tikzpicture ldots Arith s,t
%  LocalWords:  newpart forall subseteq Rightarrow forall Rightarrow subseteq t,s lf
